\documentclass[11pt]{article}
\usepackage[a4paper,margin=28mm]{geometry}
\usepackage{amsmath,amssymb,amsthm,mathtools}
\usepackage[hidelinks]{hyperref}
\usepackage{enumitem}
\newtheorem{theorem}{Theorem}
\newtheorem{lemma}{Lemma}
\newtheorem{corollary}{Corollary}
\newcommand{\E}{\mathbb E}
\newcommand{\dd}{\,\mathrm d}
\title{Sharp translation monotonicity for the polygamma\texorpdfstring{\,--\,}{--}logarithmic-mean difference}
\author{Principia Math In Progress}
\date{11 September 2026}
\begin{document}
\maketitle

\begin{abstract}
We prove the two necessary and sufficient parameter ranges in
Conjecture~5.1 of Qi's 2009 survey (also Conjecture~5.1 of its 2014
published version). A Laplace-transform difference has exactly one
change of sign. Comparing its vanishing power moment against the
polygamma kernel proves both sufficient ranges, including their
endpoints. Small-argument and large-argument asymptotics prove necessity.
The sufficient inequalities also hold for arbitrary nonconstant bounded
positive random variables translated by a common parameter.
\end{abstract}

\section*{Task 1: the source problem}
For $a,b>0$ with $a\ne b$, use Qi's parameter convention
\[
L_p(a,b)=
\begin{cases}
\displaystyle\left(\frac{b^{p+1}-a^{p+1}}{(p+1)(b-a)}\right)^{1/p},
 &p\ne-1,0,\\[1ex]
\displaystyle\frac{b-a}{\log b-\log a},&p=-1,\\[1ex]
\displaystyle\exp\left(\frac{b\log b-a\log a}{b-a}-1\right),&p=0.
\end{cases}
\]
This is the generalized logarithmic mean, despite the word ``Lehmer''
in the MathDB title. In particular, it is the power mean of a uniform
random variable on $[a,b]$. Let $i\ge0$ be an integer and put
$h_i(y)=(-1)^i\psi^{(i)}(y)$, where $\psi^{(0)}=\psi$.

\begin{theorem}\label{thm:main}
Fix distinct real $s,t$. On $x>-\min\{s,t\}$, define
\[
f_{i,p,s,t}(x)=h_i\bigl(L_p(x+s,x+t)\bigr)
 -\frac1{t-s}\int_s^t h_i(x+u)\dd u.
\]
Then $f_{i,p,s,t}$ is increasing if and only if $p\le-i-1$,
and decreasing if and only if $p\ge-i$.
In the respective ranges its derivative is strictly positive and
strictly negative everywhere. Thus the same characterization applies
whether ``increasing'' and ``decreasing'' are taken strictly or weakly.
\end{theorem}

The source locations are \cite[Conjecture 5.1, p.~16]{qi2009} and
\cite[Conjecture 5.1, p.~162]{qi2014}. The conjecture concerns translation
in $x$, with $s,t$ held fixed. We prove necessity for every such fixed
pair, not just uniformly over all pairs.

\section*{A comparison lemma}
\begin{lemma}\label{lem:crossing}
Let $X$ be a nonconstant random variable with $0<a\le X\le b<\infty$.
For real $p<1$, define
\[
m=\begin{cases}(\E X^p)^{1/p},&p\ne0,\\
\exp(\E\log X),&p=0,\end{cases}
\qquad K=m^{1-p}\E X^{p-1}.
\]
Then $K>1$, and
\[
D(u)=\E e^{-uX}-K e^{-um}
\]
has exactly one zero $u_0$ on $(0,\infty)$, is negative before $u_0$,
and positive after $u_0$. Moreover,
\begin{equation}\label{eq:zero-moment}
\int_0^\infty u^{-p}D(u)\dd u=0.
\end{equation}
If $w$ is strictly increasing and the relevant integrals converge
absolutely, then
\[
\int_0^\infty u^{-p}D(u)w(u)\dd u>0.
\]
The inequality reverses if $w$ is strictly decreasing.
\end{lemma}
\begin{proof}
For $p<0$ or $0<p<1$, apply strict Jensen to $Y=X^p$ and
the strictly convex function $y^{(p-1)/p}$:
\[
\E X^{p-1}>(\E X^p)^{(p-1)/p}=m^{p-1}.
\]
For $p=0$, strict Jensen gives
$\E X^{-1}>\exp(-\E\log X)=m^{-1}$. Thus $K>1$.

The function $G(u)=e^{um}D(u)=\E e^{(m-X)u}-K$ is strictly
convex, since
\[
G''(u)=\E\bigl((m-X)^2e^{(m-X)u}\bigr)>0.
\]
Here $G(0)=1-K<0$. Since the power mean lies strictly between the
essential infimum and essential supremum of $X$, there are
$\epsilon,\delta>0$ with $\Pr(X\le m-\epsilon)\ge\delta$.
Consequently $G(u)\ge\delta e^{\epsilon u}-K\to\infty$.
Strict convexity now gives precisely one positive zero, with the
asserted signs.

Since $1-p>0$, integrating the two nonnegative terms separately gives
\[
\int_0^\infty u^{-p}D(u)\dd u
=\Gamma(1-p)\bigl(\E X^{p-1}-Km^{p-1}\bigr)=0.
\]
Subtract $w(u_0)$ times this identity. For increasing $w$, the product
$D(u)(w(u)-w(u_0))$ is strictly positive except at $u_0$.
For decreasing $w$ it is strictly negative. This proves the comparison.
\end{proof}

\section*{Proof of the sufficient ranges}
By symmetry assume $s<t$ and let $U$ be uniform on $[s,t]$.
For fixed $x>-s$ set $X=x+U$ and $m(x)=L_p(x+s,x+t)$.
Differentiating the power-mean identity, also at $p=0$, gives
\begin{equation}\label{eq:mean-derivative}
m'(x)=m(x)^{1-p}\E X^{p-1}=K.
\end{equation}
The standard positive series and Laplace integral are
\begin{equation}\label{eq:polygamma}
h_i'(y)=(i+1)!\sum_{n=0}^\infty(y+n)^{-i-2}
 =\int_0^\infty\frac{u^{i+1}}{1-e^{-u}}e^{-yu}\dd u.
\end{equation}
They hold for every $i\ge0$ and $y>0$; see, for example,
\cite[equation (1.2)]{yang2016}. Hence
\begin{equation}\label{eq:derivative}
f'(x)=Kh_i'(m)-\E h_i'(X)
 =-\int_0^\infty u^{-p}D(u)w_{i,p}(u)\dd u,
\qquad
w_{i,p}(u)=\frac{u^{i+1+p}}{1-e^{-u}}.
\end{equation}
All integrals converge absolutely: near zero the last integrand is
$O(u^i)$, and at infinity $a=x+s>0$ supplies exponential decay.

Put $\alpha=i+1+p$. A direct logarithmic differentiation yields
\[
\frac{w_{i,p}'(u)}{w_{i,p}(u)}
 =\frac\alpha u-\frac1{e^u-1}.
\]
If $p\le-i-1$, then $\alpha\le0$ and this is strictly negative.
Lemma~\ref{lem:crossing} and \eqref{eq:derivative} imply $f'(x)>0$.
If $-i\le p<1$, then $\alpha\ge1$, and $e^u-1>u$ shows that
the logarithmic derivative is strictly positive. The same lemma gives
$f'(x)<0$. This includes $p=0$ when appropriate.

For completeness, suppose $p\ge1$. Then $m\ge\E X$ and $K\le1$.
For $p>1$, the second inequality follows from Jensen applied to
$y^{(p-1)/p}$; for $p=1$, $K=1$.
By \eqref{eq:polygamma}, $h_i'$ is positive, strictly decreasing,
and strictly convex. Consequently
\[
Kh_i'(m)\le h_i'(m)\le h_i'(\E X)<\E h_i'(X),
\]
which again gives $f'(x)<0$.

\section*{Necessity for every fixed interval}
Translate $x$ so that $a=x+s\downarrow0$ is the left endpoint,
and put $d=t-s>0$, $z=a+d/2$. All parameters $d,p,i$ are fixed below.
Expansion of the defining integral of the mean about its midpoint gives
\begin{equation}\label{eq:mean-asymptotic}
L_p(a,a+d)=z+\frac{(p-1)d^2}{24z}+O(z^{-3}),
\qquad z\to\infty.
\end{equation}
Indeed $\E(z+V)^p=z^p(1+p(p-1)d^2/(24z^2)+O(z^{-4}))$
for $V$ uniform on $[-d/2,d/2]$; at $p=0$ expand
$\E\log(z+V)$ instead.
The integral test in \eqref{eq:polygamma}, also after differentiation,
gives
\[
h_i'(z)=i!z^{-i-1}+O(z^{-i-2}),\qquad
h_i''(z)=-(i+1)!z^{-i-2}+O(z^{-i-3}).
\]
Taylor expansion, with the higher derivative bounds from the same
series, therefore proves
\begin{equation}\label{eq:f-infinity}
f_{i,p,s,t}(x)=\frac{i!(i+p)d^2}{24z^{i+2}}
                  +O(z^{-i-3}).
\end{equation}
In particular $f(x)\to0$ for every real $p$.
If $p<-i$, this expansion is negative for sufficiently large $x$.
A decreasing function with limit zero must be nonnegative everywhere.
Thus decreasing monotonicity requires $p\ge-i$.

An increasing function with limit zero must be nonpositive everywhere.
We show that $p>-i-1$ instead makes $f$ positive near $a=0$.
First suppose $i\ge1$. The recurrence for the polygamma function gives
$h_i(y)=-i!y^{-i-1}+O(1)$ as $y\downarrow0$, with a bounded
remainder on every fixed bounded positive interval. On integration,
\begin{equation}\label{eq:average-small}
\E h_i(X)=-\frac{(i-1)!}{d}a^{-i}+O(1).
\end{equation}
For $-i-1<p<-1$ the mean satisfies
\[
m\sim\bigl((-p-1)d\bigr)^{-1/p}a^{1+1/p}.
\]
The inequality $p>-i-1$ implies $(i+1)(1+1/p)<i$, so
$h_i(m)=o(a^{-i})$. For $p=-1$, the exact logarithmic-mean formula
gives $h_i(m)=O((\log(1/a))^{i+1})=o(a^{-i})$.
For $p>-1$, the mean tends to a strictly positive constant, so $h_i(m)$
stays bounded. Equation~\eqref{eq:average-small} proves in all cases
\[
f(x)\sim\frac{(i-1)!}{d}a^{-i}>0.
\]

Finally, for $i=0$ the excluded range is $p>-1$. Here
\[
m\longrightarrow
\begin{cases}d(p+1)^{-1/p},&p\ne0,\\d/e,&p=0,
\end{cases}
\qquad
\E\psi(X)=\frac{\log\Gamma(a+d)-\log\Gamma(a)}d
            =\frac{\log a}{d}+O(1).
\]
Thus $f(x)=-(\log a)/d+O(1)\to+\infty$, which also excludes increasing
monotonicity. This completes the proof of Theorem~\ref{thm:main}.

\begin{corollary}\label{cor:weighted}
Let $U$ be any nonconstant bounded real random variable. Define
$m_p(x)$ as the power mean of $x+U$, including its geometric value
at $p=0$. On $x>-\operatorname*{ess\,inf}U$, the function
\[
h_i(m_p(x))-\E h_i(x+U)
\]
has strictly positive derivative if $p\le-i-1$, and strictly negative
derivative if $p\ge-i$.
\end{corollary}
\begin{proof}
The sufficient-range proof uses only boundedness, positivity, and
nonconstancy. Its differentiation and integral interchange are valid
locally on the stated domain by uniform positive lower bounds on $x+U$.
The uniform distribution was needed only for the endpoint necessity
argument above.
\end{proof}

\section*{Follow-up tasks}
These tasks follow the full solution above. They are proposed directions,
not additional completed problems or independently screened open claims.
\begin{enumerate}[label=\textbf{Task \arabic*:},start=2,leftmargin=*]
\item Determine the parameters for which $f$ or $-f$ is completely
monotonic, as suggested in Qi's Remark~5.8. Start with the two boundary
parameters $p=-i$ and $p=-i-1$.
\item For $-i-1<p<-i$, determine the number of zeros and stationary
points of $f$. In particular, prove or disprove uniqueness of its zero
and of its global minimum, and obtain bounds for their locations.
\item Determine the sharp monotonicity ranges when the uniform law is
replaced by a beta density on $[s,t]$. Identify how endpoint behavior
affects necessity beyond the sufficient ranges in
Corollary~\ref{cor:weighted}.
\item Replace $\psi^{(i)}$ by the $q$-polygamma function for $0<q<1$,
and determine the translation-monotonicity ranges for the same means
$L_p$. Account for the discrete Laplace measure and the limit $q\uparrow1$.
\end{enumerate}

\section*{Status and verification}
The live MathDB record and a dated literature audit accompany this file.
No matching prior claim was located in the sources recorded in that audit;
this is a bounded literature check, not a guarantee about all public work.
The proof above is analytic. The accompanying script checks derivative
signs and compares separately computed expressions at prescribed inputs;
those computations are supplementary and do not replace the proof.

\begin{thebibliography}{9}
\bibitem{qi2009} F. Qi,
\emph{Bounds for the ratio of two gamma functions---From Gautschi's and
Kershaw's inequalities to completely monotonic functions},
arXiv:0904.1049v1 (2009), Conjecture~5.1.
\url{https://arxiv.org/abs/0904.1049}.
\bibitem{qi2014} F. Qi,
\emph{Bounds for the Ratio of Two Gamma Functions: from Gautschi's and
Kershaw's Inequalities to Complete Monotonicity},
Turkish Journal of Analysis and Number Theory \textbf{2}(5) (2014),
152--164. \url{https://doi.org/10.12691/tjant-2-5-1}.
\bibitem{yang2016} Z.-H. Yang and S.-Z. Zheng,
\emph{Monotonicity of a mean related to polygamma functions with an
application}, Journal of Inequalities and Applications (2016), 216.
\url{https://doi.org/10.1186/s13660-016-1155-4}.
\end{thebibliography}
\end{document}
